\subsection{Nim y misère Nim}\label{rec:nim}

Varias pilas; en cada turno se quita cualquier cantidad positiva de una
sola pila. En misère, el que toma la última piedra pierde.

\tablafunc{
\codigo{nim_win(piles)} & ¿gana el primero? & $O(n)$ \\
\codigo{misere(piles)} & versión misère & $O(n)$ \\
\codigo{nim_move(piles)} & \texttt{\{pila, nuevo tamaño\}} que deja XOR 0 & $O(n)$ \\
}

\copiar{el bloque \texttt{COPIA NIM}.}
\lstinputlisting{codigo/nim.cpp}

\ojo{En misère, si todas las pilas son 0 o 1, gana el primero cuando la
cantidad de pilas es par. En otro caso se juega como Nim normal.}